% Source: https://artofproblemsolving.com/community/c3h1595684_2016_mathcounts_state_target_6
% Stable identity: aops-topic-1595684-q01
There are $\dbinom{20}{2}=190$ equally likely pairs of cards. Count the complement: sums of at least $10$. With values from $2$ to $6$, these are
\begin{itemize}
\item $4+6$: $4\cdot4=16$ pairs;
\item $5+5$: $\dbinom42=6$ pairs;
\item $5+6$: $4\cdot4=16$ pairs;
\item $6+6$: $\dbinom42=6$ pairs.
\end{itemize}
That is $16+6+16+6=44$ pairs. Hence
\[
P(\text{sum}<10)=1-\dfrac{44}{190}=\dfrac{146}{190}=\boxed{\dfrac{73}{95}}.
\]
\textbf{Where the posted work goes wrong.} For two \emph{different} values, such as $4$ and $6$, the probability is $2\cdot\dfrac{4}{20}\cdot\dfrac{4}{19}=\dfrac{32}{380}$: after drawing a $4$, all four $6$s are still in the deck. The posted work used $\dfrac{4}{20}\cdot\dfrac{3}{19}$, which is right only for two equal values. With the correct values, the complement is $\dfrac{32+12+32+12}{380}=\dfrac{88}{380}=\dfrac{22}{95}$, not $\dfrac{12}{95}$.
